\section{Introduction}

The standard metric for highway investment prioritization is the volume-to-capacity ratio.
V/C measures how close a corridor operates to its design capacity under peak-hour demand.
The ATRI bottleneck rankings translate V/C stress into dollar cost using an average truck
delay figure (\$150--225 per truck-hour), producing an estimate of the annual economic
cost of recurring congestion. This framework is well-suited to the problem it was designed
for: identifying where additional lane-miles would reduce chronic peak-hour delay.

It is poorly suited to a different and complementary problem: identifying where corridor
closures --- complete, unplanned interruptions to traffic flow --- impose freight costs that
dwarf those of routine congestion. A congested I-285 Atlanta interchange that operates
at 0.95 V/C and delays 22,000 trucks by 15 minutes per day imposes approximately
\$1.2B in annual delay cost (\$225/hr $\times$ 0.25 hr $\times$ 22,000 trucks $\times$
365 days). But a closed Donner Pass that interrupts 8,000 trucks per day for 18 hours,
50 times per year, imposes a cost of approximately \$2.4B annually --- without ever
appearing in a V/C-based bottleneck analysis, because V/C measures peak-hour flow,
not episodic closure.

The distinction matters for investment logic. Congestion relief investments (managed
lanes, interchange reconstruction) reduce delay on corridors that are open but slow.
Resilience investments (snowsheds, drainage reconstruction, flood barriers) reduce
closure frequency on corridors that are intermittently unavailable. These are different
assets, in different places, addressing different failure modes --- and the V/C framework
systematically undercounts the economic case for resilience investment.

This paper establishes a closure cost model for the 15 highest-risk corridors in the
ROUTE corpus. The model combines: (1) closure frequency from historical incident databases;
(2) closure duration distribution from the same sources; (3) freight volume (truck AADT)
at the closure point; and (4) a rerouting-penalty computation that incorporates both the
detour distance and the B1 isolation penalty \citep{ROUTE_B1} --- the degree to which
a corridor lacks an adequate alternate. The key finding is that the corridors with the
highest closure costs are not the most congested: \textbf{Donner Pass costs 4.2$\times$
more per closure event than the Dallas interchange}, despite having a lower peak V/C at
the pass itself, because the Donner B1 isolation penalty (detour = 550 miles via I-40
Flagstaff) multiplies the per-hour cost of each closure.

The companion concept introduced here is \textbf{redundancy value}: the economic value
of having an adequate alternate route. Redundancy value is computed as the difference in
closure cost between the current corridor configuration (existing alternates) and a
counterfactual with an interstate-quality alternate. For Donner Pass, the redundancy
value of a hypothetical I-70W (western Colorado to Sacramento) is \$1.9B per year.
This figure directly answers the investment case question for alternate-route construction:
at what cost does building a redundant route pay for itself?

\subsection*{Contributions}

\begin{enumerate}
\item A closure cost model integrating frequency, duration, freight volume, and B1 isolation
penalty --- the first such model applied systematically to the ROUTE corpus
\item Quantification of the V/C--closure cost gap: the top-5 closure-cost corridors are
not in the top-5 of any V/C-based bottleneck ranking
\item Redundancy value computation for the top-10 high-closure corridors, establishing the
economic floor for alternate route investment justification
\item A compound investment case: corridors where hardening and redundancy investment are
additive (not substitutable) and where neither alone captures the full available benefit
\end{enumerate}
